\section{A direct smoothness-constrained weight representation}\label{sec:weights}

Fix any reference $\delta_*>0$ and put
\[
a_j=\delta_j/\delta_*>0,\qquad
\alpha_j=\log a_j,\qquad
\expit(z)=\frac{1}{1+e^{-z}}.
\]
For $u\in(0,1)$ define the \emph{coupled spectral-weight family}
\begin{equation}\label{eq:weights-u}
w_j(u)=\frac{a_ju}{1-u+a_ju},\quad
G_u(u)=\frac1r\sum_{j=1}^r w_j(u).
\end{equation}
This is not a free $r$-dimensional collection of weights: the common scalar $u$ couples all directions.

\begin{theorem}[Characterization of the PLS spectral family]\label{thm:odds}
For $w\in(0,1)^r$, the following statements are equivalent:
\begin{enumerate}
\item There exists $\lambda>0$ such that
$w_j=\lambda\delta_j/(1+\lambda\delta_j)$ for every $j$.
\item All numbers $w_j/[\delta_j(1-w_j)]$ are equal.
\item There exists $u\in(0,1)$ for which \eqref{eq:weights-u} holds.
\end{enumerate}
The associated $\lambda$ and $u$ are unique. Arbitrarily setting all $w_j=s$ is a PLS filter for an interior $s$ only if all positive $\delta_j$ are equal.
\end{theorem}
\begin{proof}
If (i) holds, direct rearrangement yields
$w_j/(1-w_j)=\lambda\delta_j$, proving (ii).
If (ii) holds, write the positive common value as $\lambda$;
solving $w_j=\lambda\delta_j(1-w_j)$ gives (i).
Under (i) let $u=\lambda\delta_*/(1+\lambda\delta_*)$.
Then $u/(1-u)=\lambda\delta_*$, and
\[
\frac{a_ju}{1-u+a_ju}
=\frac{\lambda\delta_j}{1+\lambda\delta_j},
\]
which is (iii). Conversely, under (iii), the positive number
$\lambda=u/[\delta_*(1-u)]$ gives (i).
Each inverse formula is injective, proving uniqueness.
Finally, if every $w_j=s\in(0,1)$, condition (ii) says
$s/[(1-s)\delta_j]$ is independent of $j$; hence the $\delta_j$
must coincide. The converse is immediate.
\end{proof}

\begin{theorem}[Unique direct-$S$ solution and explicit root bracket]\label{thm:root}
For every target $0<s<1$, there exists a unique $u_s\in(0,1)$
such that $G_u(u_s)=s$. Equivalently, with
\[
x=\logit(u)=\log\!\left(\frac{u}{1-u}\right),\qquad
G(x)=\frac1r\sum_{j=1}^r\expit(x+\alpha_j),
\]
the equation $G(x_s)=s$ has one solution on the real line, satisfying
\begin{equation}\label{eq:bracket}
\logit(s)-\max_j\alpha_j
\ \le x_s\le\
\logit(s)-\min_j\alpha_j.
\end{equation}
The weights $w_j(s)=\expit(x_s+\alpha_j)$ attain exact mean $s$.
The closed extensions are $w_j(0)=0$ and $w_j(1)=1$.
\end{theorem}
\begin{proof}
Each function $u\mapsto w_j(u)$ is continuous on $[0,1]$,
takes the values zero and one at the respective endpoints, and
has derivative
\[
w'_j(u)=\frac{a_j}{(1-u+a_ju)^2}>0.
\]
Thus their average is a continuous strictly increasing bijection
from $[0,1]$ to $[0,1]$, proving existence and uniqueness.
In log-odds coordinates,
\[
G'(x)=\frac1r\sum_{j=1}^r
\expit(x+\alpha_j)[1-\expit(x+\alpha_j)]>0.
\]
Set $x_-=\logit(s)-\max_j\alpha_j$.
For every $j$,
$x_-+\alpha_j\le\logit(s)$, so
$\expit(x_-+\alpha_j)\le s$, hence $G(x_-)\le s$.
Similarly, at
$x_+=\logit(s)-\min_j\alpha_j$, every term is at least $s$,
and $G(x_+)\ge s$. Strict monotonicity gives
$x_-\le x_s\le x_+$. Substitution gives the asserted weights.
The endpoint extensions follow by letting $x$ tend to
$-\infty$ and $+\infty$.
\end{proof}

\begin{remark}[The hidden penalty coordinate]\label{rem:implicit}
The scalar root is \emph{not} an entirely new parameterization:
from the proof of Theorem~\ref{thm:odds},
\begin{equation}\label{eq:implicit}
\lambda=\frac{e^{x_s}}{\delta_*},\qquad
x_s=\log\lambda+\log\delta_*.
\end{equation}
The algorithm need not evaluate the right side of \eqref{eq:implicit}
to form the trend, but a one-dimensional implicit inversion remains.
Replacing a stable log-$\lambda$ root solve by $x$ shifts its argument
by a constant. A speed or stability advantage therefore cannot follow
from the change of root coordinate alone.
\end{remark}

\begin{theorem}[Exact reconstruction for arbitrary fixed difference mean]\label{thm:reconstruction}
For fixed $y$, $\mu$ and $s\in[0,1]$, let $w(s)$ be the weights of
Theorem~\ref{thm:root}, including its endpoint extensions. Define
$c_\mu=V^\top y-\mu\Sigma^{-1}U^\top\one_r$.
Then
\begin{equation}\label{eq:direct-fit}
\widehat\tau_{s,\mu}
=y-V\diag(w(s))c_\mu
=y-V\left[
w(s)\odot(V^\top y)
-\mu\,(w(s)\oslash\sigma)\odot(U^\top\one_r)
\right].
\end{equation}
For each interior $s$, this is \emph{exactly} the unique minimizer
of \eqref{eq:objective} at $\lambda=\Lambda_d(s)$. No explicit
$\lambda$ or matrix inverse is required to evaluate \eqref{eq:direct-fit}
once the weights and SVD are available.
At $s=0$, $\widehat\tau_{0,\mu}=y$. At $s=1$,
\begin{equation}\label{eq:limit}
\widehat\tau_{1,\mu}
=(I_n-VV^\top)y+\mu V\Sigma^{-1}U^\top\one_r,
\end{equation}
the unique least-squares projection of $y$ onto
$\{\tau:D_d\tau=\mu\one_r\}$.
\end{theorem}
\begin{proof}
Project the normal equation \eqref{eq:normal} into
$\Ran(V)$ and its orthogonal complement. Write
$z=V^\top\widehat\tau$ and $z_y=V^\top y$.
Since $D_d^\top\one_r=V\Sigma U^\top\one_r$, the normal equation
in the $j$th right singular direction is
\[
(1+\lambda\sigma_j^2)z_j
=z_{y,j}+\lambda\mu\sigma_j(U^\top\one_r)_j.
\]
Solving and rearranging yields
\[
z_j=z_{y,j}-
\frac{\lambda\sigma_j^2}{1+\lambda\sigma_j^2}
\left(z_{y,j}-\frac{\mu}{\sigma_j}(U^\top\one_r)_j\right).
\]
The orthogonal complement component is unchanged from $y$,
because $D_d$ vanishes there. Theorems~\ref{thm:odds} and
\ref{thm:root} identify the coefficient as $w_j(s)$ for
$\lambda=\Lambda_d(s)$, proving \eqref{eq:direct-fit}.
At $s=0$ all weights vanish. At $s=1$ all weights equal one
and \eqref{eq:limit} follows. The latter vector is feasible:
multiplying by $D_d=U\Sigma V^\top$ gives
$D_d\widehat\tau_{1,\mu}=\mu UU^\top\one_r=\mu\one_r$.
Every other feasible $\tau$ has the form
$\widehat\tau_{1,\mu}+v$ with $v\in\Nul(D_d)$;
$\widehat\tau_{1,\mu}-y=-V c_\mu$ lies in
$\Nul(D_d)^\perp$. By the Pythagorean identity, adding a
nonzero feasible direction increases squared distance to $y$,
proving uniqueness of the projection.
\end{proof}

\begin{corollary}[One-dimensional penalized spectrum]\label{cor:r1}
If $r=n-d=1$, the unique weight is $w_1(s)=s$ for all
$s\in[0,1]$. No scalar root solve is needed:
$\widehat\tau_{s,\mu}=y-sVc_\mu$.
\end{corollary}
\begin{proof}
The average of a single weight equals the weight itself.
Apply \eqref{eq:direct-fit}.
\end{proof}
